% unidades/09-causativa-pasiva.tex
\chapter{😩 \ruby{使役受け身}{しえきうけみ} — Causativa-Pasiva}
\unidadheader{09}{使役受け身}{させられる}{Expresar que uno fue obligado a hacer algo}

\begin{tcolorbox}[vocabbox, title={\emoji{pushpin} Vocabulario}]
\begin{tabularx}{\linewidth}{llX}
\toprule
\textbf{Japonés} & \textbf{Español} & \textbf{Tipo} \\ \midrule
\vocabitem{残業する}{ざんぎょうする}{hacer horas extra}{v. suru}
\vocabitem{我慢する}{がまんする}{aguantar / soportar}{v. suru}
\vocabitem{走る}{はしる}{correr}{v. G1}
\vocabitem{謝る}{あやまる}{disculparse}{v. G1}
\vocabitem{飲む}{のむ}{beber / tomar}{v. G1}
\vocabitem{待つ}{まつ}{esperar}{v. G1}
\vocabitem{上司}{じょうし}{jefe / superior}{sustantivo}
\bottomrule
\end{tabularx}
\end{tcolorbox}

\subsection*{\emoji{speech-balloon} Frases Clave}
\frase{\ruby{上司}{じょうし}に\ruby{残業}{ざんぎょう}させられました。}{Fui obligado a hacer horas extra por mi jefe.}
\frase{\ruby{嫌}{きら}いな\ruby{野菜}{やさい}を\ruby{食}{た}べさせられた。}{Me obligaron a comer verduras que no me gustaban.}
\frase{\ruby{長}{なが}い\ruby{間}{あいだ}\ruby{待}{ま}たされた。}{Me hicieron esperar mucho tiempo.}
\frase{\ruby{謝}{あやま}らされた。}{Me obligaron a disculparme.}

\subsection*{\emoji{gear} Gramática}
\patron{Formación: causativa + られる}
  {V-causativa + られる → V-causativa-pasiva}
  {\ej{\ruby{食}{た}べる → \ruby{食}{た}べさせる → \ruby{食}{た}べさせられる}{G2: verb causativo + られる}
   \ej{\ruby{書}{か}く → \ruby{書}{か}かせる → \ruby{書}{か}かせられる（\ruby{書}{か}かされる）}{G1: forma larga o forma corta}
   \ej{する → させる → させられる}{G3: させられる}}

\patron{Forma corta G1 (cuando la causativa ya tiene 4+ moras)}
  {V-G1: vocal final u → aされる (forma abreviada)}
  {\ej{\ruby{飲}{の}む → \ruby{飲}{の}まされる（× \ruby{飲}{の}まさせられる）}{Se evita la acumulación de sílabas}
   \ej{\ruby{待}{ま}つ → \ruby{待}{ま}たされる。\ruby{走}{はし}る → \ruby{走}{はし}らされる。}{Forma corta muy común en habla}
   \ej{\ruby{書}{か}く → \ruby{書}{か}かされる。}{× \ruby{書}{か}かさせられる (demasiado largo)}}

\comparacion{させる vs させられる}
  {させる\\\small agente hace a otro hacer algo\\\small \ruby{子供}{こども}を\ruby{食}{た}べさせる\\\textit{Hago comer al niño (yo soy el que manda)}}
  {させられる\\\small sujeto es obligado a hacer\\\small \ruby{私}{わたし}は\ruby{食}{た}べさせられる\\\textit{Me obligan a comer (yo soy la víctima)}}
  {させられる combina la fuerza del causativo (obligación) con la pasiva (ser el receptor). Siempre implica involuntariedad y frecuentemente queja.}

\coloquial{させられた (queja formal)}{させられた / させられてさ... (queja casual)}{
  〜させられた es MUY común en conversación para quejarse. Añadir さ o のに al final enfatiza la queja: 残業させられたのに、ありがとうも言われなかった。}

\culturabox{\ruby{日本}{にほん}の\ruby{職場文化}{しょくばぶんか} — Cultura laboral japonesa}{
  La causativa-pasiva refleja la realidad del ambiente laboral japonés. La presión jerárquica implícita hace que los empleados «sean obligados» a quedarse tarde, disculparse, presentar, etc. Las expresiones con させられる son frecuentes en conversaciones sobre trabajo, y reflejan la cultura del \ruby{過労}{かろう} (exceso de trabajo).
}

\lectura{\ruby{新入社員}{しんにゅうしゃいん}の\ruby{愚痴}{ぐち}\\[0.4em]
  「\ruby{今日}{きょう}も\ruby{残業}{ざんぎょう}させられたよ。それだけじゃなくて、\ruby{嫌}{きら}いな\ruby{飲}{の}み\ruby{会}{かい}にも\ruby{参加}{さんか}させられたし、\ruby{上司}{じょうし}の\ruby{代}{か}わりに\ruby{謝}{あやま}らされた。\ruby{早}{はや}く\ruby{独立}{どくりつ}したい。」}
{Las quejas del nuevo empleado\\[0.4em]
  «Hoy también me obligaron a hacer horas extra. Y no solo eso: también me obligaron a ir a la reunión de bebidas que odio, y me hicieron disculparme en lugar de mi jefe. Quiero ser independiente pronto.»}

\begin{tcolorbox}[ejerciciobox, title={\emoji{pencil} Ejercicios}]
\begin{enumerate}
  \item Forma la causativa-pasiva: \ruby{飲}{の}む、\ruby{書}{か}く、する、\ruby{待}{ま}つ
  \item ¿Forma larga o corta? \ruby{走}{はし}る (G1, 4 moras en causativa)
  \item Crea una oración quejándote de algo usando させられた.
\end{enumerate}
\textbf{Respuestas:}
1) \ruby{飲}{の}まされる・\ruby{書}{か}かされる・させられる・\ruby{待}{ま}たされる \quad
2) \ruby{走}{はし}らされる (forma corta) \quad
3) (libre)
\end{tcolorbox}

\begin{tcolorbox}[kanjibox, title={\emoji{mahjong} Kanjis}]
\begin{tabularx}{\linewidth}{cllXX}
\toprule
\textbf{Kanji} & \textbf{On} & \textbf{Kun} & \textbf{Significado} & \textbf{Vocab} \\ \midrule
\kanjirow{残}{ざん}{\ruby{のこ}{る・す}}{quedar/resto}{\ruby{残業}{ざんぎょう} / \ruby{残}{のこ}る}
\kanjirow{業}{ぎょう}{—}{oficio/karma}{\ruby{残業}{ざんぎょう} / \ruby{業}{ぎょう}\ruby{務}{む}}
\kanjirow{我}{が}{\ruby{われ・わ}{}}{yo/ego}{\ruby{我慢}{がまん} / \ruby{我}{われ}}
\kanjirow{慢}{まん}{—}{exceso/orgullo}{\ruby{我慢}{がまん} / \ruby{自慢}{じまん}}
\bottomrule
\end{tabularx}
\end{tcolorbox}
